% Part: second-order-logic
% Chapter: metatheory
% Section: introduction

\documentclass[../../../include/open-logic-section]{subfiles}

\begin{document}

\olfileid{sol}{met}{int}

\olsection{مقدمه}

منطق مرتبهٔ اول چندین ویژگی مطلوب دارد. می‌دانیم که تصمیم‌پذیر نیست،
اما دست‌کم اصل‌موضوعی‌پذیر است. یعنی دستگاه‌های اثباتی‌ای برای منطق
مرتبهٔ اول وجود دارند که درست و کامل‌اند؛ به بیان دیگر، رابطه‌ای از
!!a{derivability}ی چون~$\Proves$ به دست می‌دهند که برای هر مجموعهٔ
!!{sentence}s چون~$\Gamma$ و هر !!{sentence} چون~$!A$،
$\Gamma \Entails !A$ اگر و تنها اگر~$\Gamma \Proves !A$. از این امر،
به‌ویژه، نتیجه می‌شود که اعتبارهای منطق مرتبهٔ اول !!{computably
  enumerable} هستند. تابع محاسبه‌پذیری چون~$f\colon \Nat \to
\Sent[L]$ وجود دارد که ارزش‌های~$f$ دقیقاً همهٔ !!{sentence}s معتبر
زبان~$\Lang{L}$ هستند. دلیلش آن است که می‌توان !!{derivation}s را
برشمرد و سپس آن‌هایی را که !!{derive} می‌کنند و حاصلشان یک~!!{sentence}
منفرد است، به همان !!{sentence} نگاشت. منطق مرتبهٔ دوم از منطق مرتبهٔ اول
بیانگرتر است و ازاین‌رو به‌طور کلی به‌چنگ‌آوردن اعتبارهایش پیچیده‌تر
است. در واقع نشان خواهیم داد که منطق مرتبهٔ دوم نه‌تنها تصمیم‌ناپذیر
است، بلکه اعتبارهای آن حتی !!{computably
  enumerable} نیز نیستند. پس
هیچ دستگاه اثباتی درست و کاملی برای منطق مرتبهٔ دوم نمی‌تواند وجود
داشته باشد (هرچند دستگاه‌های درست اما ناکامل در دسترس‌اند و در واقع
موضوع‌های مهم پژوهش‌اند).

منطق مرتبهٔ اول دو ویژگی دیگر نیز دارد: فشرده است (اگر هر زیرمجموعهٔ
متناهی از مجموعه‌ای چون~$\Gamma$ از !!{sentence}s ارضاشدنی باشد، خود
$\Gamma$ ارضاشدنی است) و قضیهٔ لوونهایم--اسکولم دربارهٔ آن برقرار است
(اگر~$\Gamma$ مدلی نامتناهی داشته باشد، مدلی~!!a{denumerable} نیز
دارد). هر دو نتیجه برای منطق مرتبهٔ دوم شکست می‌خورند. باز هم دلیل این
است که منطق مرتبهٔ دوم می‌تواند واقعیت‌هایی دربارهٔ اندازهٔ
!!{domain}s بیان کند که منطق مرتبهٔ اول نمی‌تواند.

\end{document}
